\documentclass[a4paper]{article}

\def\npart{II}

\def\ntitle{Logic and Set Theory}
\def\nlecturer{I.\ B.\ Leader}

\def\nterm{Lent}
\def\nyear{2018}

\ifx \nauthor\undefined
  \def\nauthor{Qiangru Kuang}
\else
\fi

\ifx \ntitle\undefined
  \def\ntitle{Template}
\else
\fi

\ifx \nauthoremail\undefined
  \def\nauthoremail{qk206@cam.ac.uk}
\else
\fi

\ifx \ndate\undefined
  \def\ndate{\today}
\else
\fi

\title{\ntitle}
\author{\nauthor}
\date{\ndate}

%\usepackage{microtype}
\usepackage{mathtools}
\usepackage{amsthm}
\usepackage{stmaryrd}%symbols used so far: \mapsfrom
\usepackage{empheq}
\usepackage{amssymb}
\let\mathbbalt\mathbb
\let\pitchforkold\pitchfork
\usepackage{unicode-math}
\let\mathbb\mathbbalt%reset to original \mathbb
\let\pitchfork\pitchforkold

\usepackage{imakeidx}
\makeindex[intoc]

%to address the problem that Latin modern doesn't have unicode support for setminus
%https://tex.stackexchange.com/a/55205/26707
\AtBeginDocument{\renewcommand*{\setminus}{\mathbin{\backslash}}}
\AtBeginDocument{\renewcommand*{\models}{\vDash}}%for \vDash is same size as \vdash but orginal \models is larger
\AtBeginDocument{\let\Re\relax}
\AtBeginDocument{\let\Im\relax}
\AtBeginDocument{\DeclareMathOperator{\Re}{Re}}
\AtBeginDocument{\DeclareMathOperator{\Im}{Im}}
\AtBeginDocument{\let\div\relax}
\AtBeginDocument{\DeclareMathOperator{\div}{div}}

\usepackage{tikz}
\usetikzlibrary{automata,positioning}
\usepackage{pgfplots}
%some preset styles
\pgfplotsset{compat=1.15}
\pgfplotsset{centre/.append style={axis x line=middle, axis y line=middle, xlabel={$x$}, ylabel={$y$}, axis equal}}
\usepackage{tikz-cd}
\usepackage{graphicx}
\usepackage{newunicodechar}

\usepackage{fancyhdr}

\fancypagestyle{mypagestyle}{
    \fancyhf{}
    \lhead{\emph{\nouppercase{\leftmark}}}
    \rhead{}
    \cfoot{\thepage}
}
\pagestyle{mypagestyle}

\usepackage{titlesec}
\newcommand{\sectionbreak}{\clearpage} % clear page after each section
\usepackage[perpage]{footmisc}
\usepackage{blindtext}

%\reallywidehat
%https://tex.stackexchange.com/a/101136/26707
\usepackage{scalerel,stackengine}
\stackMath
\newcommand\reallywidehat[1]{%
\savestack{\tmpbox}{\stretchto{%
  \scaleto{%
    \scalerel*[\widthof{\ensuremath{#1}}]{\kern-.6pt\bigwedge\kern-.6pt}%
    {\rule[-\textheight/2]{1ex}{\textheight}}%WIDTH-LIMITED BIG WEDGE
  }{\textheight}% 
}{0.5ex}}%
\stackon[1pt]{#1}{\tmpbox}%
}

%\usepackage{braket}
\usepackage{thmtools}%restate theorem
\usepackage{hyperref}

% https://en.wikibooks.org/wiki/LaTeX/Hyperlinks
\hypersetup{
    %bookmarks=true,
    unicode=true,
    pdftitle={\ntitle},
    pdfauthor={\nauthor},
    pdfsubject={Mathematics},
    pdfcreator={\nauthor},
    pdfproducer={\nauthor},
    pdfkeywords={math maths \ntitle},
    colorlinks=true,
    linkcolor={red!50!black},
    citecolor={blue!50!black},
    urlcolor={blue!80!black}
}

\usepackage{cleveref}



% TODO: mdframed often gives bad breaks that cause empty lines. Would like to switch to tcolorbox.
% The current workaround is to set innerbottommargin=0pt.

%\usepackage[theorems]{tcolorbox}





\usepackage[framemethod=tikz]{mdframed}
\mdfdefinestyle{leftbar}{
  %nobreak=true, %dirty hack
  linewidth=1.5pt,
  linecolor=gray,
  hidealllines=true,
  leftline=true,
  leftmargin=0pt,
  innerleftmargin=5pt,
  innerrightmargin=10pt,
  innertopmargin=-5pt,
  % innerbottommargin=5pt, % original
  innerbottommargin=0pt, % temporary hack 
}
%\newmdtheoremenv[style=leftbar]{theorem}{Theorem}[section]
%\newmdtheoremenv[style=leftbar]{proposition}[theorem]{proposition}
%\newmdtheoremenv[style=leftbar]{lemma}[theorem]{Lemma}
%\newmdtheoremenv[style=leftbar]{corollary}[theorem]{corollary}

\newtheorem{theorem}{Theorem}[section]
\newtheorem{proposition}[theorem]{Proposition}
\newtheorem{lemma}[theorem]{Lemma}
\newtheorem{corollary}[theorem]{Corollary}
\newtheorem{axiom}[theorem]{Axiom}
\newtheorem*{axiom*}{Axiom}

\surroundwithmdframed[style=leftbar]{theorem}
\surroundwithmdframed[style=leftbar]{proposition}
\surroundwithmdframed[style=leftbar]{lemma}
\surroundwithmdframed[style=leftbar]{corollary}
\surroundwithmdframed[style=leftbar]{axiom}
\surroundwithmdframed[style=leftbar]{axiom*}

\theoremstyle{definition}

\newtheorem*{definition}{Definition}
\surroundwithmdframed[style=leftbar]{definition}

\newtheorem*{slogan}{Slogan}
\newtheorem*{eg}{Example}
\newtheorem*{ex}{Exercise}
\newtheorem*{remark}{Remark}
\newtheorem*{notation}{Notation}
\newtheorem*{convention}{Convention}
\newtheorem*{assumption}{Assumption}
\newtheorem*{question}{Question}
\newtheorem*{answer}{Answer}
\newtheorem*{note}{Note}
\newtheorem*{application}{Application}

%operator macros

%basic
\DeclareMathOperator{\lcm}{lcm}

%matrix
\DeclareMathOperator{\tr}{tr}
\DeclareMathOperator{\Tr}{Tr}
\DeclareMathOperator{\adj}{adj}

%algebra
\DeclareMathOperator{\Hom}{Hom}
\DeclareMathOperator{\End}{End}
\DeclareMathOperator{\id}{id}
\DeclareMathOperator{\im}{im}
\DeclarePairedDelimiter{\generation}{\langle}{\rangle}

%groups
\DeclareMathOperator{\sym}{Sym}
\DeclareMathOperator{\sgn}{sgn}
\DeclareMathOperator{\inn}{Inn}
\DeclareMathOperator{\aut}{Aut}
\DeclareMathOperator{\GL}{GL}
\DeclareMathOperator{\SL}{SL}
\DeclareMathOperator{\PGL}{PGL}
\DeclareMathOperator{\PSL}{PSL}
\DeclareMathOperator{\SU}{SU}
\DeclareMathOperator{\UU}{U}
\DeclareMathOperator{\SO}{SO}
\DeclareMathOperator{\OO}{O}
\DeclareMathOperator{\PSU}{PSU}

%hyperbolic
\DeclareMathOperator{\sech}{sech}

%field, galois heory
\DeclareMathOperator{\ch}{ch}
\DeclareMathOperator{\gal}{Gal}
\DeclareMathOperator{\emb}{Emb}



%ceiling and floor
%https://tex.stackexchange.com/a/118217/26707
\DeclarePairedDelimiter\ceil{\lceil}{\rceil}
\DeclarePairedDelimiter\floor{\lfloor}{\rfloor}


\DeclarePairedDelimiter{\innerproduct}{\langle}{\rangle}

%\DeclarePairedDelimiterX{\norm}[1]{\lVert}{\rVert}{#1}
\DeclarePairedDelimiter{\norm}{\lVert}{\rVert}



%Dirac notation
%TODO: rewrite for variable number of arguments
\DeclarePairedDelimiterX{\braket}[2]{\langle}{\rangle}{#1 \delimsize\vert #2}
\DeclarePairedDelimiterX{\braketthree}[3]{\langle}{\rangle}{#1 \delimsize\vert #2 \delimsize\vert #3}

\DeclarePairedDelimiter{\bra}{\langle}{\rvert}
\DeclarePairedDelimiter{\ket}{\lvert}{\rangle}




%macros

%general

%divide, not divide
\newcommand*{\divides}{\mid}
\newcommand*{\ndivides}{\nmid}
%vector, i.e. mathbf
%https://tex.stackexchange.com/a/45746/26707
\newcommand*{\V}[1]{{\ensuremath{\symbf{#1}}}}
%closure
\newcommand*{\cl}[1]{\overline{#1}}
%conjugate
\newcommand*{\conj}[1]{\overline{#1}}
%set complement
\newcommand*{\stcomp}[1]{\overline{#1}}
\newcommand*{\compose}{\circ}
\newcommand*{\nto}{\nrightarrow}
\newcommand*{\p}{\partial}
%embed
\newcommand*{\embed}{\hookrightarrow}
%surjection
\newcommand*{\surj}{\twoheadrightarrow}
%power set
\newcommand*{\powerset}{\mathcal{P}}

%matrix
\newcommand*{\matrixring}{\mathcal{M}}

%groups
\newcommand*{\normal}{\trianglelefteq}
%rings
\newcommand*{\ideal}{\trianglelefteq}

%fields
\renewcommand*{\C}{{\mathbb{C}}}
\newcommand*{\R}{{\mathbb{R}}}
\newcommand*{\Q}{{\mathbb{Q}}}
\newcommand*{\Z}{{\mathbb{Z}}}
\newcommand*{\N}{{\mathbb{N}}}
\newcommand*{\F}{{\mathbb{F}}}
%not really but I think this belongs here
\newcommand*{\A}{{\mathbb{A}}}

%asymptotic
\newcommand*{\bigO}{O}
\newcommand*{\smallo}{o}

%probability
\newcommand*{\prob}{\mathbb{P}}
\newcommand*{\E}{\mathbb{E}}

%vector calculus
\newcommand*{\gradient}{\V \nabla}
\newcommand*{\divergence}{\gradient \cdot}
\newcommand*{\curl}{\gradient \cdot}

%logic
\newcommand*{\yields}{\vdash}
\newcommand*{\nyields}{\nvdash}

%differential geometry
\renewcommand*{\H}{\mathbb{H}}
\newcommand*{\transversal}{\pitchfork}
\renewcommand{\d}{\mathrm{d}} % exterior derivative

%number theory
\newcommand*{\legendre}[2]{\genfrac{(}{)}{}{}{#1}{#2}}%Legendre symbol


\makeindex

\begin{document}

\begin{titlepage}
  \begin{center}
    \includegraphics[width=0.6\textwidth]{logo.jpg}\par
    \vspace{1cm}
    {\scshape\huge Mathematics Tripos \par}
    \vspace{2cm}
    {\huge Part \npart \par}
    \vspace{0.6cm}
    {\Huge \bfseries \ntitle \par}
    \vspace{1.2cm}
    {\Large\nterm, \nyear \par}
    \vspace{2cm}
    
    {\large \emph{Lectures by } \par}
    \vspace{0.2cm}
    {\Large \scshape \nlecturer}
    
    \vspace{0.5cm}
    {\large \emph{Notes by }\par}
    \vspace{0.2cm}
    {\Large \scshape \href{mailto:\nauthoremail}{\nauthor}}
 \end{center}
\end{titlepage}

\tableofcontents

\section{Propositional Logic}

Let \(P\) be a set of \emph{primitive propositions}. Unless otherwise stated,
\[
  P = \{p_1, p_2, \dots \}.
\]
The \emph{language} or \emph{set of propositions} \(L = L(P)\) is defined inductively by
\begin{enumerate}
\item \( p \in P, p \in L\),
\item \(\bot \in L\) (reads ``false''),
\item if \(p, q \in L\) then \((p \implies q) \in L\).
\end{enumerate}

\begin{eg}
  \((p_1 \implies \bot)\), \(((p_1 \implies p_2) \implies (p_1 \implies p_3))\), \(((p_1 \implies \bot) \implies \bot)\) are elements of \(L\).
\end{eg}

\begin{note}\leavevmode
  \begin{enumerate}
  \item Each proposition is a finite union string of symbols from the alphabet \((, ), \implies, p_1, p_2, \dots\).
  \item ``Inductively defined'' means more precisely that we set
    \begin{align*}
      L_1 &= P \cup \{\bot\} \\
      L_{n + 1} &= L_n \cup \{(p \implies q): p, q \in L_n\}
    \end{align*}
    and then set \(L = L_1 \cup L_2 \cup \dots\). \(L_n\) can be seen as ``things born by time \(n\)''.
  \item Each proposition is built up \emph{uniquely} from (1), (2) and (3). For example, \(((p_1 \implies p_2) \implies (p_1 \implies p_3))\) came from \((p_1 \implies p_2)\) and \((p_1 \implies p_3)\).
  \end{enumerate}
\end{note}

Note that we often omit outer brackets or use different brackets for clarity.

We can now define for example, \(\neg p\) (reads ``not \(p\)'') as an abbreviation for \(p \implies \bot\), \(p \lor q\) (reads ``\(p\) or \(q\)'') for \((\neg p) \implies q\), \(p \land q\) (reads \(p\) and \(q\)'') for \(\neg (p \implies (\neg q))\).

\subsection{Semantic Entailment}

\begin{definition}[Valuation]\index{valuation}
  A \emph{valuation} is a function \(v: L \to \{0, 1\}\) such that
  \begin{enumerate}
  \item \(v(\bot) = 0\),
  \item \(v(p \implies q) = \begin{cases} 0 & \text{if } v(p) = 1, v(q) = 0 \\ 1 & \text{otherwise} \end{cases}\) for all \(p, q \in L\).
  \end{enumerate}
\end{definition}

\begin{remark}
  On \(\{0, 1\}\), we could define a constant \(\bot\) by \(\bot = 0\) and an operation \(\implies\) by
  \[
    (a \implies b) =
    \begin{cases}
      0 & \text{if } a = 1, b = 0 \\
      1 & \text{otherwise}
    \end{cases}
  \]
  Then a valuation is a function \(L \to \{0, 1\}\) that preserves the structure (\(\bot\) and \(\implies\)), i.e.\ it is a homomorphism.
\end{remark}

\begin{proposition}\leavevmode
  \begin{enumerate}
  \item If \(v\) and \(v'\) are valuations with \(v(p) = v'(p)\) for all \( \in P\), then \(v = v'\).
  \item For any \(w: P \to \{0, 1\}\), there exists a valuation \(v\) with \(v(p) = w(p)\) for all \(p \in P\).
  \end{enumerate}
\end{proposition}

In other words, a valuation is determined by its values on \(P\) and any values will do.

\begin{proof}\leavevmode
  \begin{enumerate}
  \item We have for all \(p \in L_1\), \(v(p) = v'(p)\). But if \(v(p) = v'(p)\) and \(v(q) = v'(q)\) then \(v(p \implies q) = v'(p \implies q)\) so \(v = v'\) on \(L_2\). Continue inductively, we have \(v = v'\) on \(L_n\) for all \(n\).
  \item Set \(v(p) = w(p)\) for all \(p \in P\) and \(v(\bot) = 0\). This defines \(v\) on \(L_1\). Having defined \(v\) on \(L_2\), use \(v(p \implies q) \begin{cases} 0 & \text{if } v(p) = 1, v(q) = 0 \\ 1 & \text{otherwise} \end{cases}\) to define \(v\) on \(L_{n + 1}\).
  \end{enumerate}
\end{proof}

\begin{eg}
  In a valuation given by
  \begin{align*}
    v(p_1) &= 1 \\
    v(p_2) &= 1 \\
    v(p_n) &= 0 \text{ for all } n \geq 3
  \end{align*}
  we have \(v(\underbrace{(p_1 \implies p_2)}_{1} \implies \underbrace{p_3}_{0}) = 0\).
\end{eg}

\begin{definition}[Tautology]\index{Tautology}
  \(p\) is a \emph{tautology}, written \(\models p\) if \(v(p) = 1\) for all valuations \(p\).
\end{definition}

\begin{eg}\leavevmode
  \begin{enumerate}
  \item \(p \implies (q \implies p)\). ``A true statement is implied by anything''. Alternatively, we could make a truth table
    \begin{table}[h]
      \centering
      \begin{tabular}{c|c|c|c}
        \(v(p)\) & \(v(q)\) & \(v(q \implies p)\) & \(v(p \implies (q \implies p))\) \\\hline
        1 & 1 & 1 & 1 \\
        1 & 0 & 1 & 1 \\
        0 & 1 & 0 & 1 \\
        0 & 0 & 1 & 1
      \end{tabular}
      \caption{Truth table}
    \end{table}
  \item \((\neg \neg p) \implies p\), i.e.\ \(((p \implies \bot) \implies \bot) \implies p\). ``Law of excluded middle''.
  \item \((p \implies (q \implies r)) \implies ((p \implies q) \implies (p \implies r))\). This is an example where writing down a truth table is not so desirable. Instead, this is not a tautology only if we have \(v\) with
    \begin{align*}
      v(p \implies (q\implies r)) &= 1 \\
      v((p \implies q) \implies (q \implies r)) &= 0
    \end{align*}
    so \(v(p \implies q) = 1, v(p \implies r) = 0\) whence \(v(p) = 1, v(r) = 0\), so also \(v(q) = 1\). But then \(v(q \implies r) = 0\) so \(v(p \implies (q \implies r)) = 0\). Absurd.
  \end{enumerate}
\end{eg}

\begin{definition}[Entailment]\index{entailment}
  For \(S \subseteq L, t \in L\), we say \(s\) \emph{entails} or \emph{semantically implies} \(t\), written \(S \models t\), if \(v(s) = 1\) for all \(s \in S\) then \(v(t) = 1\) for each valuation \(v\).
\end{definition}

This says whenever all of \(S\) is true, \(t\) is true as well.

\begin{eg}
  \(\{p \implies q, q \implies r\} \models (p \implies r)\). Indeed, suppose not. So have \(v\) with \(v(p \implies q) = v(q \implies r) = 1, v(p \implies r) = 0\). Then \(v(p) = 1, v(r) = 0\), whence \(v(q) = 0\) (from \(v(q \implies r) = 1\)), so \(v(p \implies q) = 0\). Absurd.
\end{eg}

\begin{definition}[Model]\index{model}
  If \(v(t) = 1\), we say \(t\) is \emph{true in \(v\)} or that \(v\) is a \emph{model} of \(t\).

  For \(S \subseteq L\), \(v\) is a \emph{model} of \(S\) if \(v(s) = 1\) for all \(s \in S\).
\end{definition}

Using this terminology, \(S \implies t\) says that every model of \(S\) is a model of \(t\).

\begin{note}
  \(\models t\) is equivalent to \(\emptyset \models t\).
\end{note}

\subsection{Syntactic Implication}

For a notion of ``proof'', we'll need axioms and deduction rules. As axioms, we'll take
\begin{enumerate}
\item \(p \implies (q \implies p)\) for all \(p, q \in L\).
\item \((p \implies (q \implies r)) \implies ((p \implies q) \implies (p \implies r))\) for all \(p, q, r \in L\).
\item \((\neg \neg p) \implies p\) for all \(p \in L\).
\end{enumerate}

\begin{note}
  We have already checked that these are all tautologies. Sometimes we say 3 axiom schemes to mean 3 infinite sets of axioms.
\end{note}

As deduction rules, we'll take just \emph{modus ponens}: from \(p\) and \((p \implies q)\) we can deduce \(q\).

\begin{definition}[Proof]\index{proof}
  For \(S \subseteq L\) and \(t \in L\), a \emph{proof} of \(t\) from \(S\) consists of a finite sequence \(t_1, \dots, t_n\) of propositions, with \(t_n = t\) such that for every \(i\), the proposition \(t_i\) is an axiom, or a member of \(S\), or there exists \(j, k < i\) with \(t_j = (t_k \implies t_i)\).

  We say \(S\) is the \emph{hypotheses} or \emph{premises} and \(t\) is the \emph{conclusion}.
\end{definition}

\begin{definition}
  If there is a proof of \(t\) from \(S\), say \(S\) \emph{proves} or \emph{syntactically implies} \(t\), written \(S \yields t\). If \(\emptyset \yields t\), say \(t\) is a \emph{theorem}, written \(\yields t\).
\end{definition}


\printindex

\iffalse
logic is the interplay of syntax and semantics
set: stuff with sets, universe of sets

Contents

1: Propositional logic
2: Well-ordering and ordinals
3: Posets and Zorn's Lemma
4: Predicate logic
5: Set theory
6: Cardinals

Reading:
Johnstone, Notes on logic and set theory
van Dalen, Logic and structure
Hainal & Hamburger, Set theory
Forster, Logic, induction and sets
\fi

\end{document}
